% Generated draft: requires mathematical and source review.
\documentclass[11pt]{article}
\usepackage[margin=25mm]{geometry}
\usepackage{fontspec}
\setmainfont{Latin Modern Roman}
\usepackage{amsmath,amssymb,mathtools,unicode-math}
\setmathfont{Latin Modern Math}
\usepackage{enumitem,needspace,xurl,hyperref}
\setlength{\parindent}{0pt}
\setlength{\parskip}{5pt}
\setlength{\emergencystretch}{3em}
\newcommand{\N}{\mathbb{N}}
\newcommand{\Z}{\mathbb{Z}}
\newcommand{\Q}{\mathbb{Q}}
\newcommand{\R}{\mathbb{R}}
\newcommand{\C}{\mathbb{C}}
\newcommand{\F}{\mathbb{F}}
\newcommand{\A}{\mathbb{A}}
\newcommand{\PP}{\mathbb{P}}
\newcommand{\E}{\mathbb{E}}
\newcommand{\eps}{\varepsilon}
\newcommand{\dd}{\,\mathrm{d}}
\begin{document}
\section*{UnsolvedMath ID: 30004765}
\textbf{Source problem code:} catalogue entry
\subsection{Definitions and mathematical statement}
For a group \(G\), its Schur multiplier is the abelian group
\[
M(G):=H_2(G,\Z),
\]
where group homology is taken with trivial integral coefficients. For an abelian group \(A\), written additively, its exponent, when finite, is
\[
\exp(A):=\min\{n\in\Z_{>0}:na=0\text{ for every }a\in A\}.
\]
The supplied record asks whether both equalities
\[
\exp\bigl(M(H(8,12))\bigr)=32,
\qquad
\exp\bigl(M(H(7,13))\bigr)=49
\]
hold. In particular, each proposed equality asserts that the indicated integer annihilates every multiplier element and is the least positive integer doing so.

The notation \(H(8,12)\) and \(H(7,13)\) is not defined in the supplied record: no presentations, constructions, or identifying conventions for these groups are provided. Consequently, these are the exact intended exponent questions, but a fully specified mathematical statement requires the source definitions of the two groups. No finiteness or other structural hypotheses on these groups are added here. The original extracted wording asserts the equalities; the supplied corrected statement presents them as questions, which is the formulation retained here.
\subsection{Short English statement}
Determine whether the Schur multipliers of the source-designated groups \(H(8,12)\) and \(H(7,13)\) have exponents \(32\) and \(49\), respectively. The supplied record labels the questions open, but does not define the two groups; their precise identification requires human source review.
\subsection{Sources}
\begin{itemize}
\item Supplied local UnsolvedMath registry record, permanent ID \(30004765\), source problem code \texttt{catalogue entry}, title \emph{Conjectural Schur-Multiplier Exponents}, in the Oberwolfach Reports Open Problems collection. The proposed year is recorded as 2021; no proposer is named.
\item Supplied report citation: \emph{Computational Group Theory (2022)}, Oberwolfach Reports, DOI 10.4171/owr/2021/38, pp.\ 2027--2087. \url{https://doi.org/10.4171/owr/2021/38}. The literature-source list labels this reference \emph{Oberwolfach Report 38/2021}; both supplied year designations are retained.
\item Supplied literature reference: \emph{Twenty-first Issue of the Kourovka Notebook}. \url{https://alglog.org/21tkt.pdf}. No specific problem number or page within this reference is supplied.
\item Literature assessment dated \texttt{2026-08-22}: the supplied curation labels the problem open and reports that the exponent claims remain questions in the available primary literature, with the stated computational obstruction that the groups are too large for the cited computations. This is supplied dated curation, not an independently verified assessment; the record itself requests independent verification of time-sensitive status and cited claims.
\item The registry assigns the statement status \texttt{corrected\_verified} and reports that the official Oberwolfach report was checked directly and that definitions were restored. These are supplied verification claims, not verification performed for this draft; the supplied record nevertheless contains no definitions of the groups \(H(8,12)\) and \(H(7,13)\).
\end{itemize}
\end{document}
